
Two core questions guide the experiments:

\textbf{RQ1:} Does a statistically significant correlation exist between n-gram contamination exposure and benchmark score inflation?

\textbf{RQ2:} Is n-gram overlap reliably detectable in standard benchmarks using established methodology?

\subsubsection{Models and Checkpoints}

\begin{table}[htbp]
\centering
\resizebox{\columnwidth}{!}{%
\begin{tabular}{lll}
\toprule
Size & Parameters & Training Steps \\
\midrule
410M & 405M & 0 - 143,000 \\
1B & 1.0B & 0 - 143,000 \\
1.4B & 1.4B & 0 - 143,000 \\
2.8B & 2.8B & 0 - 143,000 \\
6.9B & 6.9B & 0 - 143,000 \\
12B & 11.8B & 0 - 143,000 \\
\bottomrule
\end{tabular}
}
\end{table}

\subsubsection{Benchmarks}

\begin{table}[htbp]
\centering
\resizebox{\columnwidth}{!}{%
\begin{tabular}{llll}
\toprule
Benchmark & Samples & Task Type & Contamination Relevance \\
\midrule
MMLU & 14,042 & Multiple-choice QA & High (factual overlap likely) \\
ARC-Challenge & 1,172 & Science QA & Medium (reasoning focus) \\
HellaSwag & 10,042 & Sentence completion & Medium (synthetic generation) \\
WinoGrande & 1,267 & Coreference resolution & Low (minimal verbatim overlap expected) \\
\bottomrule
\end{tabular}
}
\end{table}

\subsubsection{Evaluation Protocol}

\textbf{Benchmark evaluation:} lm-evaluation-harness with standard settings (0-shot, auto batch size).

\textbf{Capability measurement:} WikiText-103 perplexity at each checkpoint.

\textbf{N-gram detection:} 50,000 document subset of The Pile (0.006\% of full corpus).

